\itemhead{复习题三 · 第11题（教材p.90）}{XJ-B1-C03-RV-P090-Y1543}
\label{XJ-B1-C03-RV-P090-Y1543}
{\footnotesize 湘教版教材p.90 / PDF第97页 / 复习题三 / 第11题}
\field{原题}{某企业生产一种精密电子仪器的固定成本为 20000 元，每生产一台仪器需增加投入 100 元，已知总收益满足函数：
\[R(x)=\begin{cases}400x-\dfrac12x^2,&0\le x\le400,\\80000,&x>400,\end{cases}\]
其中 $x$ 是仪器的月产量。
\subq{1}将利润表示为月产量的函数。
\subq{2}当月产量为何值时，企业所获利润最大？最大利润是多少元？}
\field{参考答案}{（1）$L(x)=\begin{cases}-x^2/2+300x-20000,&0\le x\le400,\\60000-100x,&x>400.\end{cases}$，实际月产量取非负整数。（2）第一段配方为 $-(x-300)^2/2+25000$；第二段小于20000。故300台时最大利润25000元。}
{\small \field{关联知识}{F07 分段函数、绝对值及分段应用（教材帮p.86、91、92、93 / PDF95、100、101、102）；F06 最大值、最小值与值域（教材帮p.99、83、84 / PDF108、92、93）}
\field{方法与易错点}{先列转折点和各段条件，再建立分段式，检查衔接、定义域与值域。 易错：跨段套公式；超额部分按全量收费；绝对值去号漏区间。}
\field{答案依据}{答案册相关题组：PDF第24、25、26、27页（册内页23、24、25、26）；按当前题库原题逐项重算，不按旧版题号机械对套。}}
